\documentclass[11pt]{article}

\usepackage[T1]{fontenc}
\usepackage[margin=1in]{geometry}
\usepackage{amsmath,amssymb,amsthm,mathtools}
\usepackage{booktabs}
\usepackage{tabularx}
\usepackage{enumitem}
\usepackage{microtype}
\usepackage{natbib}
\usepackage[hidelinks]{hyperref}

\newtheorem{proposition}{Proposition}
\newtheorem{corollary}{Corollary}
\newtheorem{lemma}{Lemma}
\newtheorem{assumption}{Assumption}
\newtheorem{definition}{Definition}

\newcommand{\E}{\mathbb{E}}
\newcommand{\Prob}{\mathbb{P}}
\newcommand{\ind}{\mathbf{1}}
\newcommand{\logit}{\operatorname{logit}}
\newcommand{\Pois}{\operatorname{Poisson}}
\newcolumntype{Y}{>{\raggedright\arraybackslash}X}

\title{Cheap Attention, Costly Commitment:\\
Cross-Route Inference in a Dynamic Matching Market}
\author{Anonymous}
\date{September 2026}

\begin{document}

\maketitle

\begin{abstract}
Digital matching markets frequently combine low-cost expressions of short-term
interest with a more selective market for long-term commitment. This paper asks
when attention in the first route changes choices in the second. Agents face a
two-period choice between an available long-term match and a short-term match
with a higher-ranked partner. Short-term approaches are generated by endogenous
offer decisions and form a Poisson signal of a latent state governing long-term
attainability. Agents know payoffs but may use an overly diagnostic likelihood,
treating short-term attention as if it were more selective than it is. We
characterize the posterior wedge, derive the exact signal histories on which
misspecification causes inefficient delay, and embed the resulting acceptance
rate in a two-sided participation game. An equilibrium always exists; a global
slope restriction gives uniqueness, while local comparative statics identify
the participation multiplier. When the restriction fails, a fold theorem gives
sufficient conditions for the disappearance of a high-liquidity equilibrium.
A route-transparency policy can therefore have both a direct decision effect
and an equilibrium participation effect. The mechanism is side-neutral:
demographic asymmetries enter as estimable primitives, not assumptions.

\medskip
\noindent\textbf{JEL codes:} C78, D83, D91, J12.\quad
\textbf{Keywords:} matching markets, misspecified beliefs, online dating,
Bayesian learning, search, attention.
\end{abstract}

\section{Introduction}

Many digital platforms make an expression of interest nearly costless. That
design expands the set of feasible contacts, but it also changes what a contact
means. A short-term approach can be profitable for a sender across a wide range
of recipients even when long-term commitment would be profitable only for a
much narrower set. The two actions then have different selectivity. A recipient
who treats their informational content as interchangeable may infer too much
about the set of long-term partners they can attain.

This paper isolates that cross-route inference problem. A receiver observes a
count of short-term approaches from highly ranked partners and chooses between
an available long-term match and a short-term match carrying a possibility of
later conversion. The count is generated by one of two latent attainability
states. A calibrated receiver uses the true Poisson likelihood. A misspecified
receiver uses a likelihood under which short-term attention is more diagnostic
of long-term attainability. The behavioral assumption is therefore explicit:
attention itself does not create a Bayesian bias. Bias appears only because the
receiver applies the wrong statistical model to that attention.

The analysis delivers four results. First, the log-odds belief wedge is affine
in the observed count. This produces a transparent cutoff above which the
misspecified posterior exceeds the calibrated posterior. Second, the dynamic
choice has a closed-form threshold in the perceived probability of conversion.
The combination of the two cutoffs identifies the exact offer counts at which a
receiver rejects a long-term match that they would prefer under calibrated
beliefs. Third, the acceptance rate enters a two-sided participation game. An
equilibrium exists, is unique under a global contraction condition, and has
strictly increasing participation and match volume along any stable interior
branch. Fourth, when strategic complementarity is stronger, a fold condition
generates a genuine liquidity threshold: a small reduction in acceptance can
eliminate the high-participation equilibrium.

The distinction from the closest theory is substantive. \citet{antler2022}
show that coarse strategic reasoning can produce overoptimistic search,
perpetual delay, and market failure. This paper does not claim those conclusions
as new. Its object is instead \emph{cross-route inference}: the observed signal
is an endogenous count of cheap approaches, the misspecification concerns the
likelihood mapping that count into long-term attainability, and platform
transparency changes both acceptance and two-sided participation. This delivers
history-by-history welfare losses and a route-specific policy comparative
static that are absent from a model of coarse depth of reasoning.

The paper deliberately does not assume that one demographic group sends and
another receives, that attractiveness has a universal scalar ordering, or that
short-term interaction has negligible cost. The roles ``proposer'' and
``receiver'' are market roles. Sex, orientation, age, and platform-specific
differences can shift priors, arrival rates, payoffs, and costs, but the theory
does not determine the signs of those differences.

The model also draws a boundary around its welfare claim. It does not treat
disappointment as deadweight loss. Individual welfare falls only on histories
where the action chosen under the subjective likelihood has lower payoff under
the true data-generating process. ``Collapse'' refers only to the disappearance
of a high-liquidity equilibrium under the explicit fold conditions below.

\section{Related literature}

The paper connects matching theory, statistical signaling, and equilibrium
under misspecification. The distinction between observable actions and the
private information behind them is central to signaling theory
\citep{spence1973}. Here, however, a contact is primarily a statistical signal
of the receiver's market state rather than a costly separating signal chosen to
reveal the sender's type. When short-term contact is cheap and broadly issued,
it is naturally pooling rather than separating.

The behavioral friction is related to cursed equilibrium, in which agents
underappreciate the relationship between others' information and actions
\citep{eyster2005}, and to Berk--Nash equilibrium, which gives a learning
foundation for optimal behavior under a misspecified model
\citep{esponda2016}. The baseline below is a finite subjective Bayesian model;
the repeated-learning interpretation is discussed in Section~\ref{sec:extensions}.

Classical matching work establishes stability and sorting under explicit
preference and surplus conditions \citep{gale1962,becker1973}. Search models of
marriage show how frictions and reservation strategies alter sorting
\citep{burdett1997}; \citet{chade2017} survey the broader search-and-matching
literature. The two papers closest in environment are \citet{antler2022}, where
coarse reasoning produces overoptimistic search, and \citet{antler2026}, where
partners rationally learn compatibility through dating.

\begin{table}[ht]
\centering
\caption{Closest literature and contribution boundary}
\label{tab:closest-literature}
\footnotesize
\begin{tabularx}{\textwidth}{@{}>{\raggedright\arraybackslash}p{0.19\textwidth}YY@{}}
\toprule
Study & Central object & Distinction of this paper \\
\midrule
\citet{antler2022} & Coarse strategic reasoning in two-sided search; excessive
search and possible market failure. & Endogenous attention counts, a
route-specific likelihood error, exact distorted histories, and transparency.\\
\citet{antler2026} & Correct Bayesian learning about match-specific
compatibility while dating. & Misspecified inference from one route about
attainability in another route.\\
\citet{lee2015} & Scarce ``rose'' messages as deliberate preference signals. &
Abundant approaches are statistical signals generated by cutoff offer choices,
not separating messages.\\
\citet{fong2024} & Experimental effects of perceived market size and competition
on participation and selectivity. & A mechanism linking attention information
to two-sided participation and a liquidity threshold.\\
\citet{andrew2025} & Identification of subjective marriage-market beliefs from
dynamic hypothetical choices. & A structural source of belief error and its
equilibrium feedback, rather than a belief-elicitation method.\\
\bottomrule
\end{tabularx}
\end{table}

Empirical work using online-dating interactions finds structured preferences,
sorting, and aspirational pursuit \citep{hitsch2010,bruch2018}. Field evidence
also shows that information about the number of agents on each side changes
participation and selectivity \citep{fong2024}, while a job-fair experiment
documents distorted beliefs and changed search behavior on both sides of a
matching market \citep{abebe2025}. These findings motivate a model in which an
approach, an inference, participation, and an eventual match are distinct
objects.

Reference-dependent preferences could amplify post-match dissatisfaction
\citep{koszegi2006}, but they are not needed for the baseline mechanism and are
kept outside the main model. This separation matters empirically: distorted
beliefs change choices before matching, whereas reference dependence changes
experienced utility after expectations have formed.

\section{Environment}

\subsection{Agents, dates, and routes}

There are two dates, $t\in\{1,2\}$, and two sides of a matching market. We
focus on the decision of a representative receiver. The receiver can form a
long-term match with an average-ranked proposer, denoted $A$, or accept a
one-period short-term match with a higher-ranked proposer, denoted $H$. These
labels refer to the receiver's payoff ranking, not to immutable biological
types.

At date 1 the receiver faces a stylized menu. Accepting the $A$ long-term offer
creates a match in both periods. Choosing the $H$ short-term route produces a
date-1 payoff and may convert to a long-term $H$ match at date 2. If conversion
does not occur, an $A$ long-term opportunity is available at date 2 with
probability $r\in[0,1]$. The menu is best interpreted as the decision problem
conditional on encountering both routes; Section~\ref{sec:empirics} describes
how route availability can be estimated rather than normalized.

Let $L_A>0$ be the receiver's net per-period value from the $A$ long-term match,
$L_H>L_A$ the date-2 value of a converted $H$ match, and $S_H\in\mathbb{R}$ the
net date-1 value of the short-term interaction. These values include time,
health, emotional, search, and opportunity costs. The discount factor is
$\delta\in(0,1]$. Thus the value of accepting $A$ at date 1 is
\begin{equation}
V_L=(1+\delta)L_A. \label{eq:VL}
\end{equation}

If the probability of conversion is $q$, the value of the short-term route is
\begin{equation}
V_S(q)=S_H+\delta\left[qL_H+(1-q)rL_A\right]. \label{eq:VS}
\end{equation}

\subsection{Latent attainability and attention}

The receiver is uncertain about a route-specific market state
$d\in\{0,1\}$. The state summarizes how likely highly ranked short-term
partners are to find a long-term match with this receiver mutually acceptable.
The prior is
\[
\Prob(d=1)=p\in(0,1).
\]
This is uncertainty about equilibrium long-term attainability, not uncertainty
about one's directly observed physical attributes or intrinsic worth.

During a fixed exposure window, a Poisson mass $m_H>0$ of potential $H$
proposers encounters the receiver. Each draws a private cost $k$ from continuous
cdf $K_S$ and obtains state-dependent expected gross value $v_d^S$ from making
a short-term approach. A proposer sends the approach if and only if
$k\leq v_d^S$. Independent thinning of a Poisson process therefore gives
\begin{equation}
N\mid d\sim\Pois(\lambda_d),\qquad
\lambda_d=m_HK_S(v_d^S). \label{eq:true-poisson}
\end{equation}
We assume $0<K_S(v_0^S)<K_S(v_1^S)$, hence
$0<\lambda_0<\lambda_1$. Thus the count distribution is induced by optimal
sender cutoffs rather than imposed directly. When $v_1^S-v_0^S$ is small, or
$K_S$ is locally flat, the true likelihood ratio $\lambda_1/\lambda_0$ is close
to one: cheap short-term attention is weakly selective even though offer costs
are positive.
The calibrated posterior is
\begin{equation}
\pi(n)
=\Prob(d=1\mid N=n)
=\frac{p e^{-\lambda_1}\lambda_1^n}
{p e^{-\lambda_1}\lambda_1^n+(1-p)e^{-\lambda_0}\lambda_0^n}. \label{eq:true-posterior}
\end{equation}
Its log odds satisfy
\begin{equation}
\logit\pi(n)
=\logit p+n\log\frac{\lambda_1}{\lambda_0}
-(\lambda_1-\lambda_0). \label{eq:true-logodds}
\end{equation}

Conversion occurs with probability $\gamma_d$ in state $d$, where
\begin{equation}
0\leq\gamma_0<\gamma_1\leq1. \label{eq:gammas}
\end{equation}
The objective conversion probability conditional on the observed history is
\begin{equation}
\gamma(n)=\gamma_0+(\gamma_1-\gamma_0)\pi(n). \label{eq:true-gamma}
\end{equation}
Because $\lambda_1>\lambda_0$, both $\pi(n)$ and $\gamma(n)$ are strictly
increasing in $n$.

\subsection{Cross-route misspecification}

The receiver knows payoffs and the state-dependent conversion probabilities but
may use the wrong likelihood for the offer count. Specifically, the subjective
model is
\begin{equation}
N\mid d\overset{\mathrm{subj}}{\sim}\Pois(\widetilde\lambda_d),
\qquad 0<\widetilde\lambda_0<\widetilde\lambda_1. \label{eq:subjective-poisson}
\end{equation}
The intended interpretation is that the receiver treats short-term attention as
if it had been generated by a more selective, commitment-oriented process. The
subjective posterior and conversion probability are
\begin{align}
\widetilde\pi(n)
&=\frac{p e^{-\widetilde\lambda_1}\widetilde\lambda_1^n}
{p e^{-\widetilde\lambda_1}\widetilde\lambda_1^n
+(1-p)e^{-\widetilde\lambda_0}\widetilde\lambda_0^n}, \label{eq:subjective-posterior}\\
\widetilde\gamma(n)
&=\gamma_0+(\gamma_1-\gamma_0)\widetilde\pi(n). \label{eq:subjective-gamma}
\end{align}

\begin{assumption}[Overdiagnostic subjective model]\label{ass:overdiagnostic}
The subjective likelihood ratio is more responsive to the count than the true
likelihood ratio:
\[
\frac{\widetilde\lambda_1}{\widetilde\lambda_0}
>
\frac{\lambda_1}{\lambda_0}.
\]
\end{assumption}

Assumption~\ref{ass:overdiagnostic} does not imply that the subjective posterior
is above the true posterior at every count. Low counts can be interpreted too
pessimistically. The next result identifies the crossing exactly.

\begin{proposition}[Posterior wedge]\label{prop:wedge}
Under Assumption~\ref{ass:overdiagnostic}, define
\begin{align*}
\omega
&=\log\frac{\widetilde\lambda_1/\widetilde\lambda_0}
{\lambda_1/\lambda_0}>0,\\
\eta
&=(\widetilde\lambda_1-\widetilde\lambda_0)
-(\lambda_1-\lambda_0).
\end{align*}
Then
\begin{equation}
\logit\widetilde\pi(n)-\logit\pi(n)=n\omega-\eta. \label{eq:wedge}
\end{equation}
Consequently, $\widetilde\pi(n)>\pi(n)$ and
$\widetilde\gamma(n)>\gamma(n)$ if and only if $n>\eta/\omega$.
\end{proposition}

The proposition makes the behavioral content falsifiable. High attention causes
overoptimism only when the receiver's assumed likelihood is sufficiently more
diagnostic than the true one. Under correct Bayesian updating,
$\widetilde\lambda_d=\lambda_d$ and the wedge is identically zero.

\section{Dynamic choice and equilibrium}

By \eqref{eq:VL}--\eqref{eq:VS}, the receiver prefers the short-term route when
the perceived conversion probability exceeds
\begin{equation}
\gamma^*
=\frac{(1+\delta)L_A-S_H-\delta rL_A}
{\delta(L_H-rL_A)}. \label{eq:gamma-star}
\end{equation}

\begin{assumption}[Interior choice]\label{ass:interior}
$L_H>rL_A$ and $\gamma^*\in(\gamma_0,\gamma_1)$.
\end{assumption}

The first inequality makes conversion valuable relative to returning to the
average market. The second rules out cases in which one route is optimal for
every possible belief. Let
\begin{equation}
\pi^*=\frac{\gamma^*-\gamma_0}{\gamma_1-\gamma_0}\in(0,1). \label{eq:pi-star}
\end{equation}

We use the following finite subjective-equilibrium concept.

\begin{definition}[Subjective Bayesian equilibrium]\label{def:sbe}
A subjective Bayesian equilibrium consists of a receiver strategy $a(n)$,
subjective beliefs $\widetilde\pi(n)$, and an objective outcome distribution
such that:
\begin{enumerate}[label=(\roman*)]
\item $a(n)$ maximizes \eqref{eq:VL} and \eqref{eq:VS} under
      $\widetilde\pi(n)$;
\item beliefs are generated by the subjective model
      \eqref{eq:subjective-poisson}; and
\item objective outcome frequencies are generated by the true model
      \eqref{eq:true-poisson} and the strategy $a(n)$.
\end{enumerate}
The calibrated equilibrium is the special case
$\widetilde\lambda_d=\lambda_d$.
\end{definition}

The receiver's subjective and calibrated count thresholds are
\begin{align}
\widetilde n
&=\min\{n\in\mathbb N_0:\widetilde\pi(n)\geq\pi^*\},
\label{eq:ntilde}\\
n^0
&=\min\{n\in\mathbb N_0:\pi(n)\geq\pi^*\}. \label{eq:nzero}
\end{align}
Away from knife-edge equalities, these have closed forms
\begin{align}
n^0
&=\left\lceil
\frac{\logit\pi^*-\logit p+(\lambda_1-\lambda_0)}
{\log(\lambda_1/\lambda_0)}
\right\rceil, \label{eq:nzero-closed}\\
\widetilde n
&=\left\lceil
\frac{\logit\pi^*-\logit p
+(\widetilde\lambda_1-\widetilde\lambda_0)}
{\log(\widetilde\lambda_1/\widetilde\lambda_0)}
\right\rceil. \label{eq:ntilde-closed}
\end{align}

\begin{proposition}[Cross-route distortion]\label{prop:distortion}
Under Assumption~\ref{ass:interior}, the calibrated receiver chooses $S$ if and
only if $n\geq n^0$, while the misspecified receiver chooses $S$ if and only if
$n\geq\widetilde n$. If $\widetilde n<n^0$, then exactly the histories
\begin{equation}
\mathcal D=\{\widetilde n,\widetilde n+1,\ldots,n^0-1\} \label{eq:distortion-set}
\end{equation}
produce inefficient delay. For each $n\in\mathcal D$, the receiver's material
welfare loss is
\begin{equation}
\ell(n)
=V_L-V_S(\gamma(n))
=\delta(L_H-rL_A)\left[\gamma^*-\gamma(n)\right]>0. \label{eq:individual-loss}
\end{equation}
No distortion occurs if $\widetilde n=n^0$.
\end{proposition}

The result distinguishes preference from error. A receiver may rationally value
a short-term match, and a high count may rationally predict conversion. The
model calls the choice inefficient only on histories where the receiver would
choose the long-term route under the true likelihood but chooses the short-term
route under the subjective likelihood.

Let
\begin{equation}
f(n)=(1-p)e^{-\lambda_0}\frac{\lambda_0^n}{n!}
+p e^{-\lambda_1}\frac{\lambda_1^n}{n!} \label{eq:mixture-pmf}
\end{equation}
be the unconditional count distribution. Expected receiver welfare loss is
\begin{equation}
\mathcal L
=\sum_{n\in\mathcal D}f(n)\ell(n). \label{eq:expected-loss}
\end{equation}
It is strictly positive whenever $\mathcal D$ is nonempty, because every finite
Poisson count has positive probability.

\section{Two-sided participation and liquidity}\label{sec:participation}

We now embed the receiver's decision in a two-sided participation game. Let
$n_c$ denote the count at which a receiver switches away from the long-term
route. The probability that an encountered long-term proposal is accepted is
\begin{equation}
\alpha(n_c)=\Prob(N<n_c)
=\sum_{n=0}^{n_c-1}f(n). \label{eq:acceptance-rate}
\end{equation}
Hence $\alpha^0=\alpha(n^0)$ under calibrated inference and
$\widetilde\alpha=\alpha(\widetilde n)$ under the subjective model. If
$\widetilde n<n^0$, every count has positive probability and therefore
$\widetilde\alpha<\alpha^0$.

There is a unit mass on each long-term side, proposer side $P$ and receiver side
$R$. An agent on side $g\in\{P,R\}$ draws an idiosyncratic net participation
cost $c_g$ from a continuously differentiable cdf $F_g$ on $\mathbb R$, with
density $f_g$. Allowing $F_g(0)>0$ accommodates agents with a positive baseline
value of participation and prevents zero activity from being imposed by
construction. If masses $x_P$ and $x_R$ participate, potential meetings are
\begin{equation}
M(x_P,x_R)=\mu x_Px_R,\qquad \mu\in(0,1]. \label{eq:matching-tech}
\end{equation}
A potential meeting succeeds with probability $\alpha$. A successful match
gives side $g$ benefit $b_g>0$; write $B_g=\mu b_g$. Thus the participation
cutoffs are $B_P\alpha x_R$ and $B_R\alpha x_P$.

\begin{definition}[Participation equilibrium]\label{def:participation}
Given $\alpha\in(0,1]$, a participation equilibrium is a pair of participation
masses in $[0,1]^2$ satisfying
\begin{align}
x_P&=F_P(B_P\alpha x_R), \label{eq:br-p}\\
x_R&=F_R(B_R\alpha x_P). \label{eq:br-r}
\end{align}
Successful long-term match volume is
\begin{equation}
Q(\alpha)=\mu\alpha x_Px_R. \label{eq:two-sided-volume}
\end{equation}
\end{definition}

\begin{proposition}[Existence and global uniqueness]\label{prop:existence}
For every $\alpha\in(0,1]$, a participation equilibrium exists. If the
densities are bounded by $\bar f_P$ and $\bar f_R$ and
\begin{equation}
\alpha^2B_PB_R\bar f_P\bar f_R<1, \label{eq:global-contraction}
\end{equation}
the equilibrium is unique.
\end{proposition}

Existence does not depend on a contraction or on symmetry. The uniqueness
condition has a direct interpretation: the product of the two maximal
participation responses must be smaller than one. It separates ordinary market
thinning from the strong complementarity required for multiple liquidity
equilibria.

\begin{proposition}[Participation multiplier]\label{prop:comparative-statics}
Consider an interior equilibrium at which both densities are positive. Define
\begin{align*}
z_P&=B_P\alpha x_R,&
z_R&=B_R\alpha x_P,\\
\kappa_P&=\alpha B_Pf_P(z_P),&
\kappa_R&=\alpha B_Rf_R(z_R),\\
d_P&=B_Px_Rf_P(z_P),&
d_R&=B_Rx_Pf_R(z_R).
\end{align*}
If $\kappa_P\kappa_R<1$, the equilibrium is locally unique and
\begin{align}
\frac{dx_P}{d\alpha}
&=\frac{d_P+\kappa_Pd_R}{1-\kappa_P\kappa_R}>0,
\label{eq:dxp}\\
\frac{dx_R}{d\alpha}
&=\frac{d_R+\kappa_Rd_P}{1-\kappa_P\kappa_R}>0.
\label{eq:dxr}
\end{align}
Consequently $dQ/d\alpha>0$ along the stable branch.
\end{proposition}

The direct effect of a lower acceptance rate is amplified because exit on
either side lowers the return to participation on the other. Equations
\eqref{eq:dxp}--\eqref{eq:dxr} quantify that multiplier instead of assigning it
an arbitrary reduced-form magnitude.

For a genuine tipping result, substitute \eqref{eq:br-r} into
\eqref{eq:br-p} and define
\begin{equation}
H(x;\alpha)
=F_P\!\left(B_P\alpha F_R(B_R\alpha x)\right),\qquad
G(x;\alpha)=H(x;\alpha)-x. \label{eq:composed-br}
\end{equation}

\begin{proposition}[Liquidity fold]\label{prop:fold}
Suppose $F_P$ and $F_R$ are twice continuously differentiable near
$(x^*,\alpha^*)\in(0,1)^2$ and
\begin{equation}
G(x^*;\alpha^*)=0,\quad
H_x(x^*;\alpha^*)=1,\quad
H_{xx}(x^*;\alpha^*)<0,\quad
H_\alpha(x^*;\alpha^*)>0. \label{eq:fold-conditions}
\end{equation}
Then there are neighborhoods of $x^*$ and $\alpha^*$ such that: for
$\alpha>\alpha^*$ there are exactly two nearby equilibria; for
$\alpha<\alpha^*$ there is no nearby equilibrium. The lower nearby fixed point
has $H_x>1$ and is unstable under best-response iteration; the upper fixed
point has $0\leq H_x<1$ and is locally stable.
\end{proposition}

Proposition~\ref{prop:fold} is a local saddle-node theorem, not a claim that
every market has multiple equilibria. If a separate low stable equilibrium
persists, reducing $\alpha$ through $\alpha^*$ removes the high-liquidity branch
and best-response dynamics select the remaining low-liquidity state. This is
the precise sense in which cross-route inference can tip a market.

An information policy that reveals the true route-specific arrival rates sets
$\widetilde\lambda_d=\lambda_d$. In the decision problem it eliminates
$\mathcal D$, raises acceptance from $\widetilde\alpha$ to $\alpha^0$, and
removes $\mathcal L$. In the participation game it raises activity along a
stable branch and can restore a high-liquidity equilibrium across a fold.
Whether a real label changes sender behavior or count composition is an
empirical question, so these policy claims are conditional on holding the true
offer technology fixed.

\section{Numerical illustration}

This section illustrates magnitudes; it is not a calibration. Consider the
parameters in Table~\ref{tab:parameters}. Short-term attention is only weakly
diagnostic under the true model, with arrival rates $2$ and $3$, but the
receiver behaves as if the rates were $1$ and $4$.

\begin{table}[ht]
\centering
\caption{Illustrative parameter values}
\label{tab:parameters}
\begin{tabular}{lcccccccccc}
\toprule
$p$ & $\lambda_0$ & $\lambda_1$ & $\widetilde\lambda_0$
& $\widetilde\lambda_1$ & $\gamma_0$ & $\gamma_1$ & $\delta$
& $L_A$ & $L_H$ & $S_H$ \\
\midrule
$0.30$ & $2$ & $3$ & $1$ & $4$ & $0.05$ & $0.45$ & $0.90$
& $1.00$ & $1.60$ & $1.10$ \\
\bottomrule
\end{tabular}
\begin{minipage}{0.92\textwidth}
\footnotesize\vspace{0.4em}
Note: $r=0.50$. All payoff values are net and normalized.
\end{minipage}
\end{table}

Equation~\eqref{eq:gamma-star} gives $\gamma^*=0.3535$ and
$\pi^*=0.7588$. The calibrated receiver chooses the short-term route beginning
at $n^0=8$, while the misspecified receiver does so beginning at
$\widetilde n=4$. Hence $\mathcal D=\{4,5,6,7\}$. Table~\ref{tab:posteriors}
shows the relevant conversion beliefs.

\begin{table}[ht]
\centering
\caption{Conversion beliefs around the decision threshold}
\label{tab:posteriors}
\begin{tabular}{ccccc}
\toprule
$n$ & $\pi(n)$ & $\widetilde\pi(n)$ & $\gamma(n)$
& $\widetilde\gamma(n)$ \\
\midrule
3 & 0.347 & 0.577 & 0.189 & 0.281 \\
4 & 0.444 & 0.845 & 0.228 & 0.388 \\
5 & 0.545 & 0.956 & 0.268 & 0.432 \\
6 & 0.642 & 0.989 & 0.307 & 0.445 \\
7 & 0.729 & 0.997 & 0.342 & 0.449 \\
8 & 0.802 & 0.999 & 0.371 & 0.450 \\
\bottomrule
\end{tabular}
\end{table}

The distortion set has unconditional probability $0.2015$, and expected
receiver loss is $0.0201$ in normalized payoff units. The calibrated and
subjective long-term acceptance probabilities are $0.9957$ and $0.7942$.

To illustrate the two-sided feedback, set $B_P=B_R=\mu=1$ and use the symmetric
net-cost cdf
\[
F(z)=\{1+\exp[-10(z-0.55)]\}^{-1}.
\]
Symmetric equilibria solve $x=F(\alpha x)$. The high-liquidity fold occurs at
$(x^*,\alpha^*)=(0.8639,0.8506)$. At calibrated acceptance the fixed points are
$0.0042$, $0.5882$, and $0.9870$; the first and last are stable. Starting on
the high branch, the fall from $0.9957$ to $0.7942$ crosses the fold, leaving
only the low fixed point $0.0042$. Long-term match volume falls from $0.9699$
on the calibrated high branch to $0.000014$ in the remaining subjective
equilibrium. These values demonstrate the theorem's mechanism, not an empirical
calibration or a generic prediction of collapse. The replication code computes
the fixed points and checks their residuals.

\section{Empirical content}\label{sec:empirics}

The mechanism is testable only with data that distinguish attention from
commitment. A useful dataset would record profile exposure, route or stated
intent, approaches, replies, meetings, long-term proposals, and realized
conversions. Measuring exposure is essential: raw counts confound recipient
attributes with the platform's ranking and recommendation algorithm.

The model generates five predictions.
\begin{enumerate}[leftmargin=*]
\item Holding estimated long-term type fixed, an exogenous increase in
short-term attention raises perceived long-term attainability when users are not
shown route-specific base rates.
\item The effect is concentrated around the choice cutoff: receivers far below
or above $\gamma^*$ do not change actions.
\item Revealing route labels and empirical conversion rates reduces the effect
of short-term offer counts on long-term reservation behavior.
\item A belief-induced decline in long-term acceptance reduces participation on
both long-term sides; the response becomes large where the estimated product of
best-response slopes approaches one.
\item The relevant sufficient statistic is the likelihood ratio of attention
across latent attainability states, not the raw number of approaches.
\end{enumerate}

A clean design would randomize additional profile exposure or a batch of
clearly labeled, low-commitment signals and elicit beliefs before and after the
shock. The outcome is subsequent acceptance of independently generated
long-term opportunities. A second treatment would reveal the platform's
route-specific base rates. The key estimand is the interaction between the
attention shock and transparency treatment.

Observational estimates face three major threats. First, recipients who receive
more attention may truly have higher long-term attainability. Second, platform
algorithms jointly determine exposure and partner composition. Third, stated
intent may not equal actual commitment. Randomized exposure, sender fixed
effects, experimentally varied labels, and observed conversions are therefore
more informative than cross-sectional correlations between message counts and
later outcomes.

The model is intentionally side-neutral. Empirical work may estimate
$\lambda_d$, $\widetilde\lambda_d$, $\gamma_d$, payoffs, and entry costs
separately by sex, age, orientation, or platform. A claim that the mechanism is
stronger for one group requires those estimates; it is not a theorem of the
model.

\section{Extensions and scope}\label{sec:extensions}

\paragraph{Strategic signaling.}
The baseline treats an offer as a statistical signal. A genuine signaling
extension would give proposers a private commitment cost and let them choose a
costly message or bond. A separating equilibrium would require single crossing.
If both commitment types can send the same cheap short-term offer, the offer is
pooling and the posterior over commitment remains the prior.

\paragraph{Repeated learning.}
With many episodes, agents may estimate the wrong model that best approximates
observed outcomes within a restricted subjective family. A Berk--Nash
formulation would replace the fixed pair
$(\widetilde\lambda_0,\widetilde\lambda_1)$ with the Kullback--Leibler minimizer
inside that family. Selective feedback is likely important: agents observe
conversion only for short-term offers they accept.

\paragraph{Reference dependence.}
Past attention may enter a reference point even when beliefs are calibrated.
That mechanism should be modeled as a separate state variable, following the
logic of expectations-based reference dependence, rather than inserted into the
Bayesian posterior. It can affect experienced match utility without implying
that accepting a positive-surplus match is irrational.

\paragraph{Investment.}
If agents choose costly appearance or relationship investment $x$, effective
quality can be written $q=b+g(x)$ with fixed endowment $b$, concave technology
$g$, and convex cost $K$. The private first-order condition is
$K'(x)=g'(x)\E[V_q]$. When returns are mainly ordinal, simultaneous investment
can become a positional arms race and need not change relative tiers. No
universal physical target follows from the theory.

\section{Conclusion}

Cheap short-term attention and costly long-term commitment need not carry the
same information. A calibrated Bayesian accounts for that difference, so a
large volume of attention alone cannot generate systematic overconfidence. The
mechanism in this paper requires a precise error: the receiver uses a likelihood
under which short-term approaches are more selective about long-term
attainability than they truly are.

That error has sharp consequences in the two-period model. It moves the count
threshold at which a receiver chooses the short-term route, creates an exactly
characterized set of inefficient-delay histories, and lowers long-term
acceptance. In a two-sided participation game, an equilibrium always exists,
the stable comparative statics amplify the direct acceptance effect, and an
explicit fold condition identifies when a high-liquidity equilibrium
disappears. These are conditional, falsifiable claims. Stronger statements
about demographic asymmetry, dissatisfaction, or household stability require
additional primitives and separate evidence.

\appendix

\section{Proofs}

\begin{proof}[Proof of Proposition~\ref{prop:wedge}]
Bayes' rule gives
\[
\logit\pi(n)
=\logit p+n\log(\lambda_1/\lambda_0)-(\lambda_1-\lambda_0)
\]
and
\[
\logit\widetilde\pi(n)
=\logit p+n\log(\widetilde\lambda_1/\widetilde\lambda_0)
-(\widetilde\lambda_1-\widetilde\lambda_0).
\]
Subtracting produces \eqref{eq:wedge}. Assumption~\ref{ass:overdiagnostic}
implies $\omega>0$, so the difference is positive exactly when
$n>\eta/\omega$. The logit function is strictly increasing. Hence the same
inequality orders $\widetilde\pi(n)$ and $\pi(n)$. Finally,
$\widetilde\gamma(n)-\gamma(n)
=(\gamma_1-\gamma_0)[\widetilde\pi(n)-\pi(n)]$, and
$\gamma_1>\gamma_0$ completes the proof.
\end{proof}

\begin{proof}[Proof of Proposition~\ref{prop:distortion}]
Subtracting \eqref{eq:VL} from \eqref{eq:VS} yields
\[
V_S(q)-V_L
=\delta(L_H-rL_A)(q-\gamma^*).
\]
By Assumption~\ref{ass:interior}, the coefficient on $q-\gamma^*$ is positive.
Thus a receiver chooses $S$ exactly when the relevant conversion belief is at
least $\gamma^*$. Equations \eqref{eq:true-gamma},
\eqref{eq:subjective-gamma}, and \eqref{eq:pi-star} imply that this is equivalent
to the relevant posterior being at least $\pi^*$. Both posteriors are strictly
increasing in $n$, which establishes the threshold strategies in
\eqref{eq:ntilde}--\eqref{eq:nzero}.

If $\widetilde n<n^0$, a subjective receiver chooses $S$ and a calibrated
receiver chooses $L$ exactly when $\widetilde n\leq n<n^0$. On those histories,
$\gamma(n)<\gamma^*$, so evaluating the subjective action under the true model
gives
\[
V_L-V_S(\gamma(n))
=\delta(L_H-rL_A)[\gamma^*-\gamma(n)]>0.
\]
If the thresholds coincide, the two strategies coincide for every count.
\end{proof}

\begin{proof}[Proof of Proposition~\ref{prop:existence}]
For fixed $\alpha$, define
\[
T_\alpha(x_P,x_R)
=\bigl(F_P(B_P\alpha x_R),F_R(B_R\alpha x_P)\bigr).
\]
The map is continuous and maps the nonempty compact convex set $[0,1]^2$ into
itself. Brouwer's fixed-point theorem gives existence.

For uniqueness, let $x=(x_P,x_R)$ and $y=(y_P,y_R)$ be equilibria. The bounded
density assumption and the mean-value theorem imply
\begin{align*}
|x_P-y_P|&\leq B_P\alpha\bar f_P|x_R-y_R|,\\
|x_R-y_R|&\leq B_R\alpha\bar f_R|x_P-y_P|.
\end{align*}
Combining the inequalities yields
\[
|x_P-y_P|
\leq\alpha^2B_PB_R\bar f_P\bar f_R|x_P-y_P|.
\]
Under \eqref{eq:global-contraction}, this is possible only if $x_P=y_P$;
the second inequality then implies $x_R=y_R$.
\end{proof}

\begin{proof}[Proof of Proposition~\ref{prop:comparative-statics}]
Differentiating \eqref{eq:br-p}--\eqref{eq:br-r} at an interior equilibrium
gives
\begin{align*}
dx_P&=d_P\,d\alpha+\kappa_P\,dx_R,\\
dx_R&=d_R\,d\alpha+\kappa_R\,dx_P.
\end{align*}
The Jacobian determinant is $1-\kappa_P\kappa_R>0$, so the implicit-function
theorem gives local uniqueness and solving the system gives
\eqref{eq:dxp}--\eqref{eq:dxr}. All terms in their numerators are positive.
Finally,
\[
\frac{dQ}{d\alpha}
=\mu\left[x_Px_R+\alpha x_R\frac{dx_P}{d\alpha}
+\alpha x_P\frac{dx_R}{d\alpha}\right]>0.
\]
\end{proof}

\begin{proof}[Proof of Proposition~\ref{prop:fold}]
At the candidate point, $G=G_x=0$, $G_{xx}<0$, and $G_\alpha>0$. Because
$G_{xx}\neq0$, the implicit-function theorem applied to $G_x=0$ gives a unique
continuously differentiable critical point $x_c(\alpha)$ near $x^*$.
Shrinking the neighborhood if necessary, $G_{xx}<0$ throughout, so this point
is a strict local maximum. Define $m(\alpha)=G(x_c(\alpha);\alpha)$. The envelope
identity gives
\[
m'(\alpha^*)=G_\alpha(x^*;\alpha^*)>0.
\]
Thus $m(\alpha)$ is negative immediately below $\alpha^*$ and positive
immediately above it. At $\alpha^*$, strict concavity makes $G$ negative at two
sufficiently close boundary points on either side of $x^*$. By continuity the
same is true for nearby $\alpha$. Hence, below $\alpha^*$ the local maximum is
negative and there is no nearby zero; above it, the maximum is positive and the
intermediate-value theorem gives one zero on each side. Strict concavity makes
these zeros unique in the neighborhood.

At the lower zero $G_x>0$, so $H_x>1$; at the upper zero $G_x<0$, so $H_x<1$.
Because cdf derivatives and all scale parameters are nonnegative, $H_x\geq0$.
The usual one-dimensional fixed-point criterion therefore gives instability of
the lower zero and local stability of the upper zero.
\end{proof}

\section{Verification and reproducibility}\label{sec:verification}

The replication archive contains two reproducible checks on the derivations.
First, \texttt{replicate.py} recomputes every reported number, locates all fixed
points by bracketed bisection, and asserts posterior, cutoff, fixed-point, and
stability residuals. Second, \texttt{formal/ProofAudit.lean} is checked by the
Lean kernel against Mathlib. It formalizes the dynamic choice-gap identity, the
sign of the individual welfare loss, the contraction inequality used for
uniqueness, and the product monotonicity behind match-volume comparisons. The
local fold proof remains an analytic argument in the text; the code verifies
its numerical instance but is not a formalization of the implicit-function
theorem.

\bibliographystyle{apalike}
\bibliography{references}

\end{document}
